\documentclass[11pt]{article}
\usepackage[margin=2.6cm]{geometry}
\usepackage{amsmath,amssymb,amsthm}
\usepackage{booktabs,longtable}
\newcommand{\Fp}{\mathbb F_p}
\newcommand{\Ffive}{\mathbb F_5}
\newcommand{\Fnine}{\mathbb F_9}
\newcommand{\Fthree}{\mathbb F_3}
\newcommand{\mx}{\mathsf{X}}
\newcommand{\my}{\mathsf{Y}}
\newcommand{\mz}{\mathsf{Z}}
\newcommand{\Cl}{\operatorname{Cl}}
\newcommand{\Gal}{\operatorname{Gal}}
\newcommand{\Div}{\operatorname{Div}}
\newcommand{\rk}{\operatorname{rk}}
\newcommand{\disc}{\operatorname{disc}}
\usepackage{url}
\providecommand{\nolinkurl}[1]{\url{#1}}
\newtheorem{proposition}{Proposition}
\begin{document}

\begin{center}
{\Large Worksheets W1--W4}\\[2mm]
accompanying \emph{Mild $p$-class tower groups of imaginary quadratic
fields}
\end{center}

\medskip
\noindent
These worksheets contain the worked examples referenced from
the paper as Parts W1--W4; the tables of Parts W5 and W6 are
separate documents in this directory.  References of the form Proposition~2.12,
\S4.2, Section~5, or (P.$n$) refer to statements, sections, and
numbered equations of the paper; plain parenthesised numbers refer to
equations of this document.

\section*{Part W1: The principal example --- class fields and secondary norm operators}

This part and Part~W2 record the complete data of the first
field of Section~5 of the paper (the complete $p=5$ table is
Part~W5).  We specialize to
\begin{equation}\label{eq:field}
  K=\mathbb Q(\sqrt{-2800905}),\qquad p=5.
\end{equation}
The field has discriminant
\[
  \disc(K)=-11203620
\]
and
\begin{equation}\label{eq:classgroup}
  \Cl(K)\simeq (\mathbb Z/10\mathbb Z)^3.
\end{equation}
Fix generators $g_1,g_2,g_3$ for the three cyclic factors in
\eqref{eq:classgroup}.  Their images $\overline{g_1},\overline{g_2},\overline{g_3}$ in
$\Cl(K)/5$ form a basis, and we let
\[
  \chi_1,\chi_2,\chi_3\in \bigl(\Cl(K)/5\bigr)^*
       \simeq H^1(G,\Ffive)
\]
be the dual character basis, so that for example
$\chi_1(\overline{g_1})=1$ and $\chi_1(\overline{g_2})=\chi_1(\overline{g_3})=0$.  On the torsion side we use
\begin{equation}\label{eq:torsion-basis}
  e_1=2g_1,\qquad e_2=2g_2,\qquad e_3=2g_3,
\end{equation}
which is a basis of $\Cl(K)[5]$.  These are separate choices of bases of
$\Cl(K)/5$ and $\Cl(K)[5]$; no identification between these two
$\Ffive$-vector spaces is being made.

The reconstruction theorem requires the secondary norm operators attached to
the three basis characters and their three pairwise sums, namely
\begin{equation}\label{eq:sixchars}
  \chi_1,\quad \chi_2,\quad \chi_3,\quad \chi_1{+}\chi_2,\quad \chi_1{+}\chi_3,\quad \chi_2{+}\chi_3.
\end{equation}
We first record explicitly the six corresponding unramified cyclic extensions.
For a character $x$, the vector $\boldsymbol{\eta}$ in the second column is
the integer character used by PARI on
$\Cl(K)\simeq(\mathbb Z/10\mathbb Z)^3$.  Since each cyclic factor has order $10$, the PARI vector is obtained
by multiplying the coordinates of $x$ by $10/5=2$.  Its kernel is an
index-$5$ subgroup of $\Cl(K)$ whose fixed field is $L_x$.

The last column lists generators of this subgroup directly, in the fixed
coordinates determined by $g_1,g_2,g_3$.  Thus, for
example, the first row says that the kernel belonging to $\chi_1$ is generated by
$(5,0,0)$, $(0,1,0)$, and $(0,0,1)$ in
$(\mathbb Z/10\mathbb Z)^3$.
\begin{center}
\begin{tabular}{c c c}
\toprule
$x$ & PARI vector $\boldsymbol{\eta}$ & generators of the class-field kernel \\
\midrule
$\chi_1$ & $[2,0,0]$ & $(5,0,0),\ (0,1,0),\ (0,0,1)$ \\
$\chi_2$ & $[0,2,0]$ & $(1,0,0),\ (0,5,0),\ (0,0,1)$ \\
$\chi_3$ & $[0,0,2]$ & $(1,0,0),\ (0,1,0),\ (0,0,5)$ \\
$\chi_1{+}\chi_2$ & $[2,2,0]$ & $(5,0,0),\ (4,1,0),\ (0,0,1)$ \\
$\chi_1{+}\chi_3$ & $[2,0,2]$ & $(5,0,0),\ (0,1,0),\ (4,0,1)$ \\
$\chi_2{+}\chi_3$ & $[0,2,2]$ & $(1,0,0),\ (0,5,0),\ (0,4,1)$ \\
\bottomrule
\end{tabular}
\end{center}
Each displayed subgroup has index $5$, so its class field has degree $5$ over
$K$.  In all six cases the finite relative discriminant is trivial.  Hence the
extensions $L_x| K$ are unramified at every finite prime; since $K$ is
imaginary quadratic, they are everywhere unramified.

For the Ahlqvist--Carlson formula we need more than the field $L_x$: the
generator $\sigma_x\in\Gal(L_x| K)$ must be normalized by the actual character
$x$.  We therefore also record the normalization data used in the computation,
in the notation of \S4.2 of the paper: $\sigma_{\rm Artin}$
denotes the cyclic generator of $\Gal(L_x| K)$ chosen by the class-field
computation, and the vector
$\alpha=(\alpha_1,\alpha_2,\alpha_3)$ is defined by
\[
  \operatorname{Artin}(g_i)
  =\sigma_{\rm Artin}^{\alpha_i}.
\]
The scalar $\lambda$ is characterized by
$\alpha_i=\lambda x(\overline{g_i})$ in $\Ffive$; the scalar $u$ compares the
generator initially used by the implementation with $\sigma_{\rm Artin}$,
and $\nu_x$ is defined by
$\sigma_x=\sigma_{\rm Artin}^{\nu_x}$.  Thus $\nu_x$ is the exponent of the
finally normalized generator.  In the present computation $u=1$ in every
case, so $\nu_x=\lambda$.
\begin{center}
\begin{tabular}{c c c c c}
\toprule
$x$ & $\alpha$ & $\lambda$ & $u$ & $\nu_x$ \\
\midrule
$\chi_1$   & $[3,0,0]$ & $3$ & $1$ & $3$ \\
$\chi_2$   & $[0,4,0]$ & $4$ & $1$ & $4$ \\
$\chi_3$   & $[0,0,4]$ & $4$ & $1$ & $4$ \\
$\chi_1{+}\chi_2$ & $[3,3,0]$ & $3$ & $1$ & $3$ \\
$\chi_1{+}\chi_3$ & $[3,0,3]$ & $3$ & $1$ & $3$ \\
$\chi_2{+}\chi_3$ & $[0,2,2]$ & $2$ & $1$ & $2$ \\
\bottomrule
\end{tabular}
\end{center}
As a direct check of the normalization, for every row and every
$i=1,2,3$ we verified
\[
  \operatorname{Artin}(g_i)
  =\sigma_x^{x(\overline{g_i})}.
\]

\begin{proposition}[Secondary norm matrices]\label{prop:Dmatrices}
Represent an element of $\Cl(K)/5$ by its coordinate column vector with
respect to $\overline{g_1},\overline{g_2},\overline{g_3}$.  With respect to the
input basis $e_1,e_2,e_3$ of $\Cl(K)[5]$, the six secondary norm operators
$D_x\colon \Cl(K)[5]\to\Cl(K)/5$ are represented by the matrices
\begin{align*}
D_{\chi_1}&=
\begin{pmatrix}0&0&0\\0&3&1\\0&3&0\end{pmatrix},&
D_{\chi_2}&=
\begin{pmatrix}0&1&2\\0&0&0\\4&2&4\end{pmatrix},&
D_{\chi_3}&=
\begin{pmatrix}1&2&0\\4&0&1\\0&0&0\end{pmatrix},\\[2mm]
D_{\chi_1+\chi_2}&=
\begin{pmatrix}0&3&1\\0&2&4\\0&1&2\end{pmatrix},&
D_{\chi_1+\chi_3}&=
\begin{pmatrix}1&4&0\\1&4&4\\4&1&0\end{pmatrix},&
D_{\chi_2+\chi_3}&=
\begin{pmatrix}3&1&2\\0&3&2\\0&2&3\end{pmatrix}.
\end{align*}
\end{proposition}

\begin{proof}
For each of the $18$ pairs $(x,e_j)$, the constructed pair
$(t_{x,j},I'_{x,j})$ satisfies the two exact equations
(P.46)--(P.47).  By Proposition~2.12 of the paper, these verified
equations imply
\[
  D_x(e_j)=[N_x(I'_{x,j})]\in\Cl(K)/5.
\]
The norm ideal on the right is recomputed independently and its class is
expressed in the coordinates
$\overline{g_1},\overline{g_2},\overline{g_3}$.  The resulting values are
\begin{center}
\begin{tabular}{c c c c}
\toprule
$x$ & $D_x(e_1)$ & $D_x(e_2)$ & $D_x(e_3)$ \\
\midrule
$\chi_1$   & $[0,0,0]$ & $[0,3,3]$ & $[0,1,0]$ \\
$\chi_2$   & $[0,0,4]$ & $[1,0,2]$ & $[2,0,4]$ \\
$\chi_3$   & $[1,4,0]$ & $[2,0,0]$ & $[0,1,0]$ \\
$\chi_1{+}\chi_2$ & $[0,0,0]$ & $[3,2,1]$ & $[1,4,2]$ \\
$\chi_1{+}\chi_3$ & $[1,1,4]$ & $[4,4,1]$ & $[0,4,0]$ \\
$\chi_2{+}\chi_3$ & $[3,0,0]$ & $[1,3,2]$ & $[2,2,3]$ \\
\bottomrule
\end{tabular}
\end{center}
For example, the second entry in the row for $\chi_1$ says
\[
  D_{\chi_1}(e_2)=3\overline{g_2}+3\overline{g_3}\in\Cl(K)/5,
\]
and this is the second column of the displayed matrix $D_{\chi_1}$.  Reading the
three values in each row as coordinate columns gives the six matrices
above.  The matrices were assembled only after all norm classes had been
computed and independently extracted.
\end{proof}

As an additional arithmetic check, we repeated the complete class-field and
secondary-norm computation for the three doubled characters and obtained
\begin{equation}\label{eq:scaling-check}
  D_{2\chi_1}=4D_{\chi_1},\qquad D_{2\chi_2}=4D_{\chi_2},\qquad D_{2\chi_3}=4D_{\chi_3}.
\end{equation}
By the quadratic dependence (P.14) these identities are forced;
confirming them establishes nothing new, and their worth lies in the fact that
they could have failed.  The pairs $x$ and $2x$ define the same class-field
kernel but require different normalized automorphisms, and a computation that
attached the automorphism to $\ker x$ rather than to $x$ would return
$D_{2x}=D_x$ here.  The check therefore excludes one specific way for the
implementation to go wrong at its most delicate step, and no more.  What it
probes belongs to the procedure rather than to the field, so we
applied this check only to the principal example.

\section*{Part W2: The principal example --- the transversality computation}

We apply Corollary~3.15 of the paper to
\[
  x=a+b.
\]
From the fourth matrix in Proposition~\ref{prop:Dmatrices},
\[
  \rk D_{\chi_1+\chi_2}=1,
\]
with
\[
  \ker D_{\chi_1+\chi_2}
  =\langle e_1,\,3e_2+e_3\rangle
\]
and
\[
  \operatorname{im}D_{\chi_1+\chi_2}
  =\left\langle
    \overline{g_1}+4\overline{g_2}+2\overline{g_3}
  \right\rangle.
\]
Take $v_1=a$, $v_2=c$, and $y=b+3c$.  The character $y$ annihilates the
image line above, while its restriction to $(a+b)^\perp$ is nonzero.  With
the adapted datum
\[
  \mathcal B=(e_1,3e_2+e_3;a,c;b+3c),
\]
the transversality matrix is
\[
  B_{a+b,\mathcal B}=
  \begin{pmatrix}
    3&2\\
    3&4
  \end{pmatrix},
  \qquad
  \det B_{a+b,\mathcal B}=1\in\Ffive.
\]
Thus $\chi_1{+}\chi_2$ is transverse, and Theorem~3.13 of the paper proves that
$G_{\varnothing}(K)(5)$ is mild, and hence has cohomological dimension $2$.

As an independent check, the six matrices also reconstruct a cubic word
matrix of rank three.  A direct change of variables and relation basis gives
three relations with combinatorially free leading words
\[
  \mz\mz\mx,\qquad \mz\my\my,\qquad \mz\my\mx.
\]
The accompanying computation checks the shuffle identities, the word order,
and Anick's non-overlap condition directly.  These additional calculations
are not needed for the preceding proof.

\section*{Part W3: The norm-degeneracy scheme of the principal example}

The norm-degeneracy scheme of Definition~3.17 of the paper can be
written out completely for the principal example.  With $x=ua+vb+wc$
in the basis of Part~W1, polarizing the six
matrices of Proposition~\ref{prop:Dmatrices} gives the matrix of quadratic forms
\begingroup\small
\[
  D_{(u,v,w)}=
  \begin{pmatrix}
    2vw+w^2 & 2uv+2uw+v^2+3vw+2w^2 & 4uv+2v^2\\[1pt]
    2uw+vw+4w^2 & 3u^2+4uv+uw+3vw & u^2+3uv+2uw+vw+w^2\\[1pt]
    uv+4uw+4v^2+vw & 3u^2+uv+3uw+2v^2 & 3uv+4v^2+4vw
  \end{pmatrix},
\]
\endgroup
whose $j$th column is $D_x(e_j)$ in the coordinates
$\overline{g_1},\overline{g_2},\overline{g_3}$.  The scheme
$\Sigma_D\subseteq\mathbb P^2_{\Ffive}$ is defined by the nine
$2\times2$ minors, the homogeneous quartics
\begingroup\small
\begin{align*}
 &2u^2vw+4u^2w^2+4uv^2w+4uw^3+4v^3w+4v^2w^2+4vw^3+2w^4,\\
 &4u^2vw+u^2w^2+3uv^2w+uvw^2+2uw^3+3v^3w+4v^2w^2+3vw^3+w^4,\\
 &2u^3w+4u^2vw+u^2w^2+4uv^2w+uvw^2+uw^3+4v^2w^2+2w^4,\\
 &3u^2v^2+u^2vw+uv^3+uvw^2+v^4+v^3w+v^2w^2+3vw^3,\\
 &u^2v^2+4u^2vw+2uv^3+4uv^2w+3uvw^2+2v^4+v^3w+2v^2w^2+4vw^3,\\
 &3u^3v+u^2v^2+4u^2vw+uv^3+4uv^2w+4uvw^2+v^3w+3vw^3,\\
 &2u^3v+4u^3w+4u^2v^2+4u^2w^2+4uv^3+4uv^2w+4uvw^2+2uw^3,\\
 &4u^3v+u^3w+3u^2v^2+u^2vw+2u^2w^2+3uv^3+4uv^2w+3uvw^2+uw^3,\\
 &2u^4+4u^3v+u^3w+4u^2v^2+u^2vw+u^2w^2+4uv^2w+2uw^3.
\end{align*}
\endgroup

Its closed points can be listed explicitly.  On the chart $w=1$ the
radical of the minor ideal has a reduced coordinate ring of
$\Ffive$-dimension three, which a single closed point of degree three
exhausts; on the line $w=0$ the ideal restricts to a squarefree
quadratic with two rational roots, and $(1{:}0{:}0)$ does not lie on
$\Sigma_D$.  Thus $\Sigma_D$ has exactly three closed points: the
rational points $(1{:}1{:}0)$ and $(1{:}3{:}0)$ and one point of
degree three.  At $(1{:}1{:}0)$ --- the character $\chi_1{+}\chi_2$ of
Part~W2 --- the tangent space is zero: the
point is reduced and isolated, which is the geometric form
(Corollary~3.23 of the paper) of the determinant computation
carried out there.  At $(1{:}3{:}0)$ the operator also has rank one,
but the tangent space is one-dimensional, so this point is not
transverse: rank one alone does not suffice for mildness.  The
degree-three point is again reduced and isolated, with transversality
over its residue field $\mathbb F_{125}$; either of the two
transverse points proves Theorem~5.1 of the paper for this field.  The
display is generated deterministically from the stored matrices by
\nolinkurl{tools/sigma_d_equations.py}.

\section*{Part W4: One certificate entry, item by item}

To make the certificate format concrete, we follow a single entry from
prescribed character to certified value.  The table abbreviates
two bulky items: the ten elements of $B_L$ are described but not
listed, and the $128$ factors of $t$ are counted but not printed.
The autogenerated companion \nolinkurl{certificate-guide.pdf},
included alongside this document, restates all $18$ entries of
this field in the notation of the paper --- it, too, keeps $t$ in
compact form ---, and the data file \nolinkurl{certificate.gp} is
the certificate itself, byte for byte.  Entry~$2$ certifies
the value $D_{\chi_1}(e_2)$, that is, the second column of the first matrix
in Proposition~\ref{prop:Dmatrices}.

The rows follow the entry tuple of the certificate definition in
the paper's appendix~C; the table shows the stored data, lightly
annotated (norms, factor counts, residue degrees), and the checks
the verifier runs on them are listed after it.

\begin{center}
\begin{tabular}{@{}p{0.24\textwidth}p{0.68\textwidth}@{}}
\toprule
tuple item & stored data (entry $2$)\\
\midrule
$x$ &
the character vector $(1,0,0)$, that is, $x=\chi_1$; stored under
the label \texttt{a}\\[2pt]
$e$ &
the basis element $e_2$, stored as the column index $2$\\[2pt]
$f_{\mathrm{rel}}$ &
$y^5-y^4-1027y^3+(8s-33869)y^2+(340s-648684)y+(2688s-3346560)$
over $K=\mathbb Q(s)$, $s^2+2800905=0$\\[2pt]
$f_{\mathrm{abs}}$ &
$y^{10}-2y^9-2053y^8-65684y^7-174901y^6+64171174y^5
+2665457137y^4+66051314232y^3+1091726153376y^2
+9461326049280y+31436965969920$ over $\mathbb Q$\\[2pt]
$B_L$ &
ten elements of $\mathbb Q[\theta]$, $\theta$ the class of $y$
in the absolute model, each written as a rational number or a
polynomial in $\theta$; the basis of $K$ is stored once in the
certificate header\\[2pt]
$\sigma$ &
the coordinate vector $(0,-1,0,0,0,0,0,0,-1,0)$ with respect to
$B_L$\\[2pt]
$a'$, $J$ &
$a'=\frac{41298191}{2498705294406001}
+\frac{16828}{2498705294406001}\,s$ and
$J=\bigl(\begin{smallmatrix}1201&1072\\[1pt]0&1\end{smallmatrix}\bigr)$
in Hermite normal form, $N(J)=1201$\\[2pt]
$t$, $I'$ &
$t$ in compact form, a product of $128$ factors with signed
exponents, never expanded; $I'$ an ideal of $L_x$ of absolute
norm $10587476640=2^5\cdot3\cdot5\cdot13\cdot1696711$, in
Hermite normal form\\[2pt]
$\ell$, $\mathfrak q$ &
$\ell=11$ and a degree-one prime $\mathfrak q$ above it\\[2pt]
$v$ &
$(0,3,3)$\\
\bottomrule
\end{tabular}
\end{center}

On these data the verifier decides:
\begin{itemize}
\item the change of basis between $B_L$ and a recomputed integral
  basis of the absolute model is integral with determinant $\pm1$
  --- an integral basis is not canonical, and this check is what
  lets the stored coordinates denote the same algebraic objects
  everywhere;
\item $f_{\mathrm{abs}}$ is the absolute defining polynomial
  derived from $f_{\mathrm{rel}}$, the extension $L_x\mid K$ is
  of degree $5$ and unramified everywhere, and the exact Artin
  kernel check identifies $L_x$ as the fixed field of $\ker x$;
\item the element denoted by $\sigma$ is an automorphism of exact
  order $5$ fixing $K$, normalized to $x$ as in (P.22) through
  exact Artin exponents;
\item $(a')J^5=\mathcal O_K$ --- the condition in (P.4) written
  multiplicatively --- and $[J]$ represents $e_2$;
\item $I'^{\,(1-\sigma_x)^2}\,(t)\,J\mathcal O_{L_x}
  =\mathcal O_{L_x}$, identity (P.46), holds as an equality of
  fractional ideals, evaluated factor by factor from the compact
  $t$;
\item the principal ideal generated by $N_x(t)/a'$ equals
  $\mathcal O_K$ --- computed, not inferred from (P.46) --- and
  the sign is resolved as in Lemma~4.2 of the paper: since the
  torsion units of $K$ are $\pm1$, the reduction at
  $\mathfrak q$, which equals the residue of $+1$, proves
  $N_x(t)=a'$, identity (P.47);
\item the recomputed class $[N_x(I')]$ has coordinate vector
  $v=(0,3,3)$ in $\Cl(K)/5$; by (P.26) this is
  $D_{\chi_1}(e_2)$, the second column of $D_{\chi_1}$ in
  Proposition~\ref{prop:Dmatrices}.
\end{itemize}

Entries $1$--$18$ certify the six matrices of Proposition~\ref{prop:Dmatrices}
column by column in this way.

\section*{Parts W5 and W6: the tables}

The complete tables --- Part~W5 for $p=5$ and Part~W6 for $p=3$ ---
are the separate documents \nolinkurl{tables-p5.pdf} and
\nolinkurl{tables-p3.pdf} in this directory.  Part~W6 is generated
from the result records of the repository by
\nolinkurl{generate_tables_p3.py}; regenerating it after an
extension of the range reproduces the tables from the records
alone.

\end{document}
